% supp_math.tex -- Appendices S1 and S2: results that the main theorem does not use
% (Section 2.6, the end of Section 4.3, the end of Section 5.10 and Appendix C of the first revision).
% S1.3 is app_extra_A.tex, S2 is app_extra_B.tex.
\section{Supplements to Sections 2 and 4}\label{sec:extra}

This appendix contains three supplements to \cref{sec:board,sec:leaf}: the Computed Fact that $\G$ is the
full automorphism group of the board (\cref{ssec:aut}), a description in words of the test for $\LEAF$
that \cref{alg:leaf} states as pseudocode (\cref{ssec:supp-leafwords}), and the monotonicity of the pairing
condition $\CC$ in \Breaker's stones (\cref{sec:extra-monoC}). None of them is used in the proof of
\cref{thm:main} or of \cref{cor:seven}.

\subsection{The full automorphism group of the board}\label{ssec:aut}

An \emph{automorphism} of the board is a bijection of $\Q$ that maps every line onto a line. By
\cref{lem:group} every element of $\G$ is an automorphism. The proof of the main theorem uses nothing
more, but it is natural to ask whether $\G$ misses any symmetry.

\begin{fact}[The full automorphism group; not used in the proof of \cref{thm:main}]\label{fact:aut}
The automorphism group of the incidence structure $(\Q,\Lines)$ has order $1{,}152$; hence it is equal
to~$\G$.

\established the program \fn{autgroup.cpp} in the directory \fn{logs/canon/} of the
archive (output \fn{autgroup.log}; see \cref{ssec:artifact}), a
search by individualization and color refinement on the incidence graph with $216+148$ vertices. It
finds a base consisting of three lines and three cells; the orbit of each base point under the pointwise
stabilizer of the preceding ones has $4$, $3$, $2$, $6$, $4$, $2$ elements, respectively, so that the
group has order $4\cdot3\cdot2\cdot6\cdot4\cdot2=1{,}152$.
\end{fact}

For the $4\times4\times4$ board the automorphism group, which has order $192$, was determined by
Silver~\cite{Silver1967}, and the automorphism groups of the boards $n^d$ in general were characterized
by Dvo\v{r}\'ak and Valla~\cite{DvorakValla2021}. We did not rely on these results, and we have not
compared the order $1{,}152$ with the general description in~\cite{DvorakValla2021}:
\cref{fact:aut} is a direct computation for the present board.

\subsection{The leaf test in words}\label{ssec:supp-leafwords}

The results of \cref{ssec:deciding} (\cref{prop:triples,prop:maxflow,cor:pruning}) give the following procedure for deciding $\LEAF(B,F,k)$, where
$\theta=k-|F|\ge 0$; \cref{alg:leaf} states it as pseudocode. All its steps are computations with small integers on a bipartite graph with at
most $148+216$ vertices.
\begin{enumerate}[label=\textup{\arabic*.}]
\item Test whether $F$ is $B$-admissible (\cref{def:admissible}). If not, $\LEAF(B,F,k)$ is false.
\item Compute a maximum assignment for $(B,F)$ by augmenting paths (\cref{lem:hall}). If some live line
      receives fewer than two cells, there is no pairing and $\LEAF(B,F,k)$ is false. Otherwise a pairing
      $\pi$ has been found.
\item Optionally, apply the quick rejection of \cref{cor:pruning}\,(b) to the $F$-components of
      $\live{B}$ with at least three lines.
\item For every $F$-connected triple $T$ (a line together with two lines that are $F$-adjacent to it),
      decide whether $\mu(T)\ge\theta$: starting from $\pi$, try to perform $\theta$ augmentations from
      the lines of $T$ (\cref{prop:maxflow}\,(d) or (e)). A triple may be skipped if one of its lines or
      pairs of lines is already known to satisfy $\mu\ge\theta$ (\cref{cor:pruning}\,(a)). If the
      procedure stops after fewer than $\theta$ augmentations, then $R=R_\alpha(T)$ is an $F$-connected
      set of at least three live lines with $\sur_F(R)=\mu(T)<\theta$ (\cref{prop:maxflow}\,(c)), and
      $\LEAF(B,F,k)$ is false.
\item If every triple passes, $\LEAF(B,F,k)$ is true (\cref{prop:triples}\,(b)).
\end{enumerate}
A negative answer in steps 3 and 4 comes with a set $S$ of lines whose surplus can be recomputed from
the definition. What a positive answer comes with, and why only positive answers matter for the
soundness of the main theorem, is said at the end of \cref{ssec:deciding} and in \cref{rem:true-cert}.

% app_extra_A.tex -- Appendix S1.3 (input by supp_math.tex): supplement to Section 4, not used by the main theorem.

\subsection{Monotonicity of \texorpdfstring{$\CC$}{C} in Breaker's stones}\label{sec:extra-monoC}

The condition $\LEAF$ is monotone in the set of \Breaker's stones: $\LEAF(B,F,k)$ implies
$\LEAF(B+b,F,k)$ for every cell $b\notin F$ (\cref{prop:leaf-monotone}\,(c)). In this subsection we settle
the corresponding question for the pairing condition $\CC$ of \cref{def:C} itself. The result is not used in
the proof of \cref{thm:main}, nor anywhere else in the paper; we include it because the question is
natural and the answer is not what one might expect at first.

One is tempted to regard the implication $\CC(B,F,k)\Rightarrow\CC(B+b,F,k)$ as obvious, since an
additional stone of \Breaker{} only kills lines. The difficulty is that the family of sets $X$ over
which $\CC$ quantifies grows when $B$ grows: a set that is $(B+b)$-admissible need not be $B$-admissible.
The lines through $b$ are dead with respect to $B+b$, so $X$ may have five cells on such a line, or four
cells on it and four cells on a live line that crosses it in an empty cell; both are forbidden with
respect to $B$. The statement $\CC(B,F,k)$ says nothing about such sets $X$, and the implication does
not follow by restriction. In fact it is false in general.

\begin{proposition}[Monotonicity of $\CC$ in $B$]\label{prop:C-monotone}
Let $B,F$ be disjoint sets of cells, $k\ge|F|$, and $b$ a cell with $b\notin B\cup F$.
\begin{enumerate}[label=\textup{(\alph*)}]
\item If $F$ is $B$-admissible, then $\CC(B,F,k)$ implies $\CC(B+b,F,k)$.
\item Without the hypothesis that $F$ is $B$-admissible the implication can fail: if $l$ is a line,
      $b\in l$, $B=\emptyset$ and $F=l\setminus\{b\}$, then $\CC(\emptyset,F,5)$ holds and
      $\CC(\{b\},F,5)$ does not.
\end{enumerate}
\end{proposition}

\begin{proof}[Proof of \textup{(b)}]
With respect to $B=\emptyset$ the line $l$ is live and contains the five cells of $F$; so $F$ is not
$\emptyset$-admissible, no set $X\supseteq F$ is (\cref{lem:adm-closed}\,(c)), and $\CC(\emptyset,F,5)$ holds
vacuously. With respect to $\{b\}$ the line $l$ is dead. Every other line meets $l$ in at most one cell
(\cref{lem:geometry}\,(c)) and thus contains at most one cell of $F$; so $F$ is $\{b\}$-admissible, and
$X=F$ is one of the sets to which $\CC(\{b\},F,5)$ refers. There is, however, no pairing for $(\{b\},F)$: the cell
$b$ lies on at most seven lines, so at least $141$ lines are live and need $282$ cells, while only
$216-6=210$ cells lie outside $\{b\}\cup F$ (\cref{lem:slack-bound}).
\end{proof}

\begin{proof}[Proof of \textup{(a)}]
Assume that $F$ is $B$-admissible and that $\CC(B,F,k)$ holds. Let $X\supseteq F$ be a
$(B+b)$-admissible set with $|X|\le k$; in particular $b\notin X$ and $X\cap B=\emptyset$. We have to show
that there is a pairing for $(B+b,X)$, that is, by \cref{lem:hall}\,(b), that $\dfc_X(S)\le 0$ for every
set $S\subseteq\live{B+b}$. Fix such a set $S$. We shall construct a set $X^\circ$ of cells and a set
$S'$ of lines such that
\begin{equation}\label{eq:monoC-goal}
  F\subseteq X^\circ\subseteq X,\qquad X^\circ\text{ is $B$-admissible},\qquad S'\subseteq\live{B},\qquad
  \dfc_{X^\circ}(S')\ge\dfc_X(S).
\end{equation}
This suffices: since $|X^\circ|\le k$, the statement $\CC(B,F,k)$ provides a pairing for $(B,X^\circ)$, so
$\dfc_{X^\circ}(S')\le 0$ by \cref{lem:hall}\,(b), and hence $\dfc_X(S)\le 0$.

\smallskip\noindent\emph{Notation.}
Let $D$ be the set of the lines through $b$ that are live with respect to $B$. Then $\live{B}$ is the
disjoint union of $\live{B+b}$ and $D$, and any two lines of $D$ meet in $b$ only
(\cref{lem:geometry}\,(c)). Every line $l\in D$ has at most five cells in $X$, because $b\in l\setminus X$.
Put
\[
  D_5=\{l\in D : |l\cap X|=5\},\qquad D_4=\{l\in D : |l\cap X|=4\},
\]
and for $l\in D_4$ let $c_l$ denote the cell with $l\setminus X=\{b,c_l\}$. We shall use two facts.
\begin{enumerate}
\item[\textup{($\alpha$)}] Every line of $D$ has at most four cells in $F$, and at most one line of $D$ has exactly four.
      Indeed, the lines of $D$ are live with respect to $B$ and $F$ is $B$-admissible; two lines of $D$
      with four cells in $F$ each would be a danger pair of $F$ with common cell $b\notin F$.
\item[\textup{($\beta$)}] For $l\in D_4$ there is at most one line $a\in\live{B+b}$ that passes through $c_l$ and has
      $|a\cap X|=4$. Indeed, two such lines would be a danger pair of $X$ with respect to $B+b$, with
      common cell $c_l\notin X$. If this line exists, we denote it by $a_l$.
\end{enumerate}

\smallskip\noindent\emph{The trivial case.}
If $D_4\cup D_5=\emptyset$, we take $X^\circ=X$ and $S'=S$. The set $X$ is $B$-admissible: the lines of
$\live{B+b}$ have at most four cells in $X$, the lines of $D$ at most three, and a danger pair of $X$
with respect to $B$ would consist of two lines of $\live{B+b}$ and would be a danger pair with respect
to $B+b$. So \eqref{eq:monoC-goal} holds.

\smallskip\noindent\emph{Construction of $X^\circ$.}
Now let $D_4\cup D_5\ne\emptyset$. We choose a line $l_0\in D_4\cup D_5$ as follows: if some line of $D$
has four cells in $F$, let $l_0$ be this line (it belongs to $D_4\cup D_5$ because $F\subseteq X$);
otherwise let $l_0$ be a line of $D_5$ if $D_5\ne\emptyset$, and any line of $D_4$ if $D_5=\emptyset$. By
($\alpha$),
\begin{equation}\label{eq:monoC-three}
  |F\cap l|\le 3\quad\text{for every } l\in D\setminus\{l_0\},\qquad |F\cap l_0|\le 4 .
\end{equation}
Let $R\subseteq X\setminus F$ be a set consisting of the following cells:
\begin{enumerate}[label=\textup{(C\arabic*)}]
\item two cells of $(l\cap X)\setminus F$ for every $l\in D_5\setminus\{l_0\}$, and one cell of
      $(l\cap X)\setminus F$ for every $l\in D_4\setminus\{l_0\}$;
\item if $l_0\in D_5$: one cell $r_0$ of $(l_0\cap X)\setminus F$;
\item if $l_0\in D_4$ and the line $a_0:=a_{l_0}$ of \textup{($\beta$)} exists: one cell $r^*$ of
      $\bigl((l_0\cup a_0)\cap X\bigr)\setminus F$.
\end{enumerate}
These cells exist. For (C1) and (C2) this follows from \eqref{eq:monoC-three}. For (C3), suppose that
$(l_0\cup a_0)\cap X\subseteq F$. Then $F$ contains the four cells of $l_0\cap X$ and the four cells of
$a_0\cap X$, and no further cell of these lines, as $F\subseteq X$. Both lines are live with respect to
$B$, and they share the cell $c_{l_0}\notin X$. So they form a danger pair of $F$ with respect to $B$,
contradicting the hypothesis on $F$.

Let $\varepsilon=1$ if case (C3) applies and $\varepsilon=0$ otherwise, and put $X^\circ=X\setminus R$. Then
$F\subseteq X^\circ\subseteq X$, and in all cases
\begin{equation}\label{eq:monoC-R}
  |R|\le 2|D_5|+|D_4|-1+\varepsilon .
\end{equation}
(If $l_0\in D_5$, then $R$ has at most $2(|D_5|-1)+|D_4|+1$ cells and $\varepsilon=0$; if $l_0\in D_4$, it
has at most $2|D_5|+(|D_4|-1)+\varepsilon$ cells.)

\smallskip\noindent\emph{Claim 1: $X^\circ$ is $B$-admissible.}
First, $X^\circ\cap B=\emptyset$. Next, every line that is live with respect to $B$ has at most four cells
in $X^\circ$: the lines of $\live{B+b}$ have at most four cells in $X\supseteq X^\circ$; a line of
$D\setminus\{l_0\}$ has at most three cells in $X^\circ$, by (C1) if it belongs to $D_4\cup D_5$ and
because it has at most three cells in $X$ otherwise; and $l_0$ has at most four cells in $X^\circ$, by
(C2) if $l_0\in D_5$. In particular
\begin{equation}\label{eq:monoC-l0}
  \text{$l_0$ is the only line of $D$ that can have four cells in $X^\circ$.}
\end{equation}
We first note a property of the lines of $\live{B+b}$: if $a\in\live{B+b}$ has exactly four cells in
$X^\circ$, then it has at least four, hence exactly four cells in $X$, so that
\begin{equation}\label{eq:monoC-m}
  a\cap X=a\cap X^\circ,\qquad a\cap R=\emptyset,
\end{equation}
and a cell of $a$ that is not in $X^\circ$ is not in $X$.

Suppose now that $m,m'$ is a danger pair of $X^\circ$ with respect to $B$, with common cell
$c\notin X^\circ$. By \eqref{eq:monoC-l0} the lines $m,m'$ do not both belong to $D$. If both belong to
$\live{B+b}$, then by \eqref{eq:monoC-m} each has exactly four cells in $X$, and $c\notin X$; so $m,m'$ is a
danger pair of $X$ with respect to $B+b$, which is impossible. Hence, after renaming, $m=l_0$ and
$m'\in\live{B+b}$; the line $l_0$ has exactly four cells in $X^\circ$, and \eqref{eq:monoC-m} holds for
$a=m'$. Moreover $c\ne b$, because $b\notin m'$.
\begin{itemize}
\item If $l_0\in D_5$, then $l_0\setminus X^\circ=\{b,r_0\}$, since $r_0$ is the only cell of $R$ on $l_0$
      (the cells chosen in (C1) lie on other lines of $D$, which meet $l_0$ in $b$ only). Hence
      $c=r_0\in m'\cap R$, contradicting \eqref{eq:monoC-m}.
\item If $l_0\in D_4$, then no cell of $l_0$ belongs to $R$, because $l_0$ has four cells in $X$ and four
      in $X^\circ$. Hence $l_0\setminus X^\circ=\{b,c_{l_0}\}$ and $c=c_{l_0}$. The line $m'\in\live{B+b}$
      passes through $c_{l_0}$ and has exactly four cells in $X$, so $m'=a_0$ by ($\beta$); in particular
      $a_0$ exists and $\varepsilon=1$. Then, however, the cell $r^*\in R$ lies on $l_0$ or on $a_0=m'$,
      contradicting either the first sentence of this item or \eqref{eq:monoC-m}.
\end{itemize}
This proves Claim 1.

\smallskip\noindent\emph{Claim 2: the set
$S'=S\cup D_4\cup D_5$, enlarged by the line $a_0$ if $\varepsilon=1$, satisfies
$\dfc_{X^\circ}(S')\ge\dfc_X(S)$.}
We have $S'\subseteq\live{B}$, since $\live{B}=\live{B+b}\cup D$. By the definition of the deficiency the
claim is equivalent to
\begin{equation}\label{eq:monoC-claim2}
  |\N(S')\setminus X^\circ|-|\N(S)\setminus X|\ \le\ 2\,(|S'|-|S|) .
\end{equation}
A cell of $\N(S')\setminus X^\circ$ lies in $R$ or in $\N(S')\setminus X$, and
\[
  \N(S')\setminus X=(\N(S)\setminus X)\ \cup\ \bigl(\{b\}\cup\{c_l : l\in D_4\}\bigr)\ \cup\ (a_0\setminus X),
\]
where the last set is present only if $\varepsilon=1$; in this case $a_0\setminus X=\{c_{l_0},z_0\}$ for
some cell $z_0$, because $|a_0\cap X|=4$. Recall that $S\cap D=\emptyset$, since $S\subseteq\live{B+b}$. We
use \eqref{eq:monoC-R} and distinguish three cases.
\begin{itemize}
\item $\varepsilon=0$. Then $|S'|-|S|=|D_4|+|D_5|$, and the left side of \eqref{eq:monoC-claim2} is at most
      $|R|+1+|D_4|\le(2|D_5|+|D_4|-1)+1+|D_4|=2\,(|D_4|+|D_5|)$.
\item $\varepsilon=1$ and $a_0\notin S$. Then $|S'|-|S|=|D_4|+|D_5|+1$, and the left side is at most
      $|R|+1+|D_4|+1\le(2|D_5|+|D_4|)+|D_4|+2=2\,(|D_4|+|D_5|+1)$; here the last summand $1$ accounts for
      the cell $z_0$, the cell $c_{l_0}$ being counted among the cells $c_l$.
\item $\varepsilon=1$ and $a_0\in S$. Then $|S'|-|S|=|D_4|+|D_5|$, and the cells $c_{l_0}$ and $z_0$ belong
      to $\N(S)\setminus X$ already. So the left side is at most
      $|R|+1+(|D_4|-1)\le(2|D_5|+|D_4|)+|D_4|=2\,(|D_4|+|D_5|)$.
\end{itemize}
In each case \eqref{eq:monoC-claim2} holds. Claims 1 and 2 give \eqref{eq:monoC-goal}, and the proof is
complete.
\end{proof}

In a play in which \Maker's stones $X$ are admissible at the beginning of every round, the revealed set
$F\subseteq X$ is always admissible with respect to the current set $B$ (\cref{lem:adm-closed}). So part
(a) applies in every position that the strategy of \cref{sec:game} can meet, and along such a play the
condition $\CC$, once it holds, continues to hold when \Breaker{} places further stones. The proofs of this
paper never rely on this. The only monotonicity statements about leaf conditions that occur in them
are those of \cref{prop:leaf-monotone} for $\LEAF$, whose proofs take a few lines.

As a safeguard against an oversight in the case analysis, part (a) was also tested by brute force on
small instances (boards $4^3$, $5^3$, $6^2$ and $6^3$ restricted to few live lines, with the definitions
transferred to lines of length $n$): the script \fn{logs/proof/fuzz_cprime.py} of the archive examined $5{,}574$ triples
$(B,F,k)$ for which $\CC(B,F,k)$ holds, together with $447{,}456$ sets $X$ that are $(B+b)$-admissible, of
which $244{,}718$ are not $B$-admissible, and found a pairing for $(B+b,X)$ in every case. This
computation is not part of any proof.


\section{Supplements to Section 5}\label{sec:extra-game}

This appendix contains two supplements to \cref{sec:game}: lazy monotonicity and its corollary for the
eager variant of the abstract game, with proofs (\cref{sec:extra-lazy}), and the exact number of stones
of the patterns $\Etwo$ and $\Ethree$ on the $6\times6\times6$ board (\cref{sec:extra-e23}). Neither is
used in the proof of \cref{thm:main} or of \cref{cor:seven}.

% app_extra_B.tex -- Appendix S2 (input by supp_math.tex): supplements to Section 5, not used by the main theorem.

\subsection{Lazy monotonicity and the eager variant}\label{sec:extra-lazy}

The abstract game is monotone in the revealed set: revealing more stones never hurts \Breaker. This is not
needed for \cref{thm:main}, but it clarifies the relation between lazy revelation and variants in
which more information is used. The proofs of the next two statements are given at the end of this
subsection.

\begin{proposition}[Lazy monotonicity]\label{prop:lazy-monotone}
Fix $\Rmax$. Let $F\subseteq F'$, where $F'$ is $B$-admissible and $|F'|\le t$.
\begin{enumerate}[label=\textup{(\alph*)},leftmargin=2.4em]
\item If the P-state $\Pst(B,F,t)$ is \WON, then $\Pst(B,F',t)$ is \WON.
\item If the Q-state $\Qst(B,F,t)$ is \WON, then $\Qst(B,F',t)$ is \WON.
\end{enumerate}
\end{proposition}

In an \emph{eager} variant of the abstract game, a stone that \Maker{} places on a cell of a fixed set
$Z\subseteq\Q$ is revealed at once; the natural choice for $Z$ is the set $\Busy$ of the $24$ busy cells. Formally, let the label $\WONp$ be defined like \WON{} in
\cref{def:won}, with the same $\Rmax$ and with the following additional requirement in (P2):
\begin{enumerate}[leftmargin=2.6em]
\item[\textup{(b)}] for every cell $h\in Z\setminus(B\cup F)$, the Q-state $\Qst(B,F+h,t+1)$ is $\WONp$ or
  \VOID.
\end{enumerate}
The recursion is well founded by the rank of \cref{lem:rank}, since $\Qst(B,F+h,t+1)$ has smaller rank
than $\Pst(B,F,t)$.

\begin{corollary}[Eager variant]\label{cor:eager}
For every set $Z\subseteq\Q$, every state that is \WON{} is $\WONp$. In particular, if the root is \WON{} for
$\Rmax$, then it is $\WONp$ for $\Rmax$.
\end{corollary}

\begin{remark}[Lazy versus eager]\label{rem:eager}
The label $\WONp$ keeps all requirements of \WON{} and adds the children \textup{(b)}. In a game in which
every stone of \Maker{} on $Z$ is really revealed at once, requirement \textup{(ii)} would be vacuous for table
moves $m\in Z$, because an unrevealed stone never lies on $Z$; so $\WONp$ is at least as demanding as such a
game. \cref{cor:eager} thus says that a lazy certificate is the stronger statement. The converse is
not claimed, and the eager variant was not investigated computationally. In terms of strategies, eager
revelation is the special case of lazy revelation in which \Breaker{} inspects all cells of $Z$ after every
move of \Maker: at any time \Breaker{} may replace $F$ in the record by a set $F'$ with
$F\subseteq F'\subseteq X$. The statements \ref{I:B}--\ref{I:adm} of \cref{thm:soundness} are not
affected, and $\Pst(B,F',t)$ is still \WON{} by \cref{prop:lazy-monotone}; but its canonical form
need not belong to the reachable set of the given certificate, so a certificate for the eager variant has to
provide additional table entries, with up to $|Z|$ additional children at every P-state. Lazy revelation
avoids this branching and treats table moves on busy and on non-busy cells alike.
\cref{thm:soundness} does not depend on \cref{prop:lazy-monotone}.
\end{remark}

\medskip\noindent\textbf{Proofs.}
We now prove \cref{prop:lazy-monotone} and
\cref{cor:eager}. The notation is that of \cref{sec:game}; the label \WON{} refers to a fixed
$\Rmax$, and $\rk$ is the rank of \cref{lem:rank}. Neither statement is used in the proof of
\cref{thm:main}.

\begin{lemma}[$\LOST$ is antitone in $F$]\label{lem:lost-antitone}
Let $F\subseteq F'$ be sets of cells disjoint from $B$. If $\LOST(B,F',t)$ holds, then $\LOST(B,F,t)$ holds.
\end{lemma}

\begin{proof}
Let a pattern for $(B,F')$ fit $(B,F',t)$. The same lines, cells and stone sets form a pattern for $(B,F)$:
the conditions (E2.1), (E2.2), (E3.1), (E3.2) of \cref{def:lost} do not involve the revealed set,
and the conditions (E2.3), (E3.3), which are of the form ``$e\notin F'$'' or ``$F'\cap a\subseteq S_a$'', remain
true when $F'$ is replaced by a subset. The union of $F$ and the stone sets is contained in the union of
$F'$ and the stone sets, so it has at most $t+1$ cells.
\end{proof}

\begin{proof}[Proof of \cref{prop:lazy-monotone}]
The set $F$ is $B$-admissible, because $F\subseteq F'$ and $F'$ is $B$-admissible
(\cref{lem:adm-closed}\,(a)). Hence none of the four states that occur in (a) and (b) is \VOID. We prove (a)
and (b) simultaneously by induction on the rank of the state with the smaller set $F$, that is, of
$\Pst(B,F,t)$ in (a) and of $\Qst(B,F,t)$ in (b).

(a) Let $\Pst(B,F,t)$ be \WON. If $\LEAF(B,F,t)$ holds, then $\LEAF(B,F',t)$ holds by the monotonicity of
$\LEAF$ in the revealed set (\cref{prop:leaf-monotone}\,(a), which applies because $F'$ is
$B$-admissible and $|F'|\le t$), and $\Pst(B,F',t)$ satisfies (P1). Otherwise $\Pst(B,F,t)$ satisfies (P2):
$t<\Rmax$, $\LOST(B,F,t)$ is false, and every child of $\Pst(B,F,t)$ is \WON{} or \VOID. We verify (P2) for
$\Pst(B,F',t)$. First, $\LOST(B,F',t)$ is false by \cref{lem:lost-antitone}.

\emph{Child \textup{(a)}.} The child $\Qst(B,F,t+1)$ of $\Pst(B,F,t)$ is not \VOID, hence \WON. It has
smaller rank than $\Pst(B,F,t)$, and $|F'|\le t+1$. By induction, part (b), $\Qst(B,F',t+1)$ is \WON.

\emph{Children \textup{(c)}.} Let $l\in\live{B}$ and $e\in l\setminus F'$ with
$|F'\cup(l\setminus\{e\})|\le t+1$ define a child \textup{(c)} of $\Pst(B,F',t)$, namely
$\Pst(B+e,F_1',t+1)$ with $F_1'=F'\cup(l\setminus\{e\})$, and suppose that this child is not \VOID, that is,
$F_1'$ is $(B+e)$-admissible. Since $e\notin F$ and
$|F\cup(l\setminus\{e\})|\le|F'\cup(l\setminus\{e\})|\le t+1$, the same data define a child \textup{(c)}
of $\Pst(B,F,t)$, namely $\Pst(B+e,F_1,t+1)$ with $F_1=F\cup(l\setminus\{e\})\subseteq F_1'$. This child is
not \VOID{} (\cref{lem:adm-closed}\,(a)), hence it is \WON, and it has smaller rank than $\Pst(B,F,t)$. As
$F_1\subseteq F_1'$, $F_1'$ is $(B+e)$-admissible and $|F_1'|\le t+1$, the state $\Pst(B+e,F_1',t+1)$ is \WON{}
by induction, part (a).

\emph{Children \textup{(d)}.} Let $a,b,c,Q_a,Q_b$ define a child \textup{(d)} of $\Pst(B,F',t)$ that is not
\VOID. From $c\notin F'$, $F'\cap a\subseteq Q_a$, $F'\cap b\subseteq Q_b$ and $|F'\cup Q_a\cup Q_b|\le t+1$ we
get the same four conditions for $F$. So the same data define a child \textup{(d)} of $\Pst(B,F,t)$, whose
revealed set $F\cup Q_a\cup Q_b$ is contained in $F'\cup Q_a\cup Q_b$. The argument is completed as for
\textup{(c)}.

Thus every child of $\Pst(B,F',t)$ is \WON{} or \VOID, and $\Pst(B,F',t)$ satisfies (P2).

(b) Let $\Qst(B,F,t)$ be \WON, and let $m\notin B\cup F$ be a table move for it.

\emph{Case $m\in F'$.} Then $m\in F'\setminus F$, so $|F|<|F'|\le t$, and the child \textup{(ii)} of
$\Qst(B,F,t)$ via $m$, the state $\Qst(B,F+m,t)$, exists. It is not \VOID, because $F+m\subseteq F'$; hence
it is \WON{} by requirement (ii) of \cref{def:won}. Its rank is smaller than that of $\Qst(B,F,t)$,
and $F+m\subseteq F'$. By induction, part (b), $\Qst(B,F',t)$ is \WON.

\emph{Case $m\notin F'$.} We show that $m$ is a table move for $\Qst(B,F',t)$; note $m\notin B\cup F'$.
Requirement (i): the state $\Pst(B+m,F,t)$ is \WON{} and has smaller rank than $\Qst(B,F,t)$; the set $F'$ is
$(B+m)$-admissible because $m\notin F'$ (\cref{lem:adm-closed}\,(b)), and $|F'|\le t$. By induction, part
(a), $\Pst(B+m,F',t)$ is \WON. Requirement (ii): suppose that $|F'|<t$ and that $\Qst(B,F'+m,t)$ is not
\VOID, that is, $F'+m$ is $B$-admissible. Then $|F|<t$, so the child \textup{(ii)} of $\Qst(B,F,t)$ via $m$,
the state $\Qst(B,F+m,t)$, exists; it is not \VOID, because $F+m\subseteq F'+m$; hence it is \WON, and its
rank is smaller than that of $\Qst(B,F,t)$. As $F+m\subseteq F'+m$ and $|F'+m|\le t$, the state
$\Qst(B,F'+m,t)$ is \WON{} by induction, part (b).
\end{proof}

\begin{proof}[Proof of \cref{cor:eager}]
We show by induction on the rank that every state that is \WON{} is $\WONp$. Let $\Pst(B,F,t)$ be \WON. If
it satisfies (P1), it is $\WONp$, since (P1) is the same for both labels. Otherwise it satisfies (P2); its
children \textup{(a)}, \textup{(c)}, \textup{(d)} are \WON{} or \VOID, hence $\WONp$ or \VOID{} by induction.
It remains to consider a child \textup{(b)}, that is, $\Qst(B,F+h,t+1)$ with $h\in Z\setminus(B\cup F)$. If
$F+h$ is not $B$-admissible, this state is \VOID. Otherwise observe that the child \textup{(a)},
$\Qst(B,F,t+1)$, is \WON, and that $F\subseteq F+h$, where $F+h$ is $B$-admissible and $|F+h|\le t+1$. By
\cref{prop:lazy-monotone}\,(b), $\Qst(B,F+h,t+1)$ is \WON. Its rank is smaller than that of
$\Pst(B,F,t)$, so it is $\WONp$ by induction. Finally let $\Qst(B,F,t)$ be \WON, with table move $m$. Its
child \textup{(i)} via $m$ is \WON, hence $\WONp$ by induction, and its child \textup{(ii)} via $m$, if it
exists, is \WON{} or \VOID, hence $\WONp$ or \VOID. So $\Qst(B,F,t)$ is $\WONp$.
\end{proof}

\subsection{The patterns \texorpdfstring{$\Etwo$}{E2} and \texorpdfstring{$\Ethree$}{E3} on the
\texorpdfstring{$6\times6\times6$}{6x6x6} board}\label{sec:extra-e23}

\cref{prop:lost-stones} bounds the number of distinct stones of a pattern from below by $11$ for
$\Etwo$ and by $12$ for $\Ethree$, using only that two lines share at most one cell. On the
$6\times6\times6$ board the first bound is sharp and the second is not. Nothing in the paper depends on
this.

For the proof we regard cells as points of $\mathbb{R}^3$. Let
$D=\{-1,0,1\}^3\setminus\{(0,0,0)\}$. Recall that a line is a set $\{(f_1(t),f_2(t),f_3(t)):0\le t\le 5\}$,
where each $f_i$ is a constant, $t$, or $5-t$, and not all three are constant. Such a line lies on a real
line with direction vector $d\in D$, where $d_i$ is $0$, $1$ or $-1$ according as $f_i$ is constant, $t$ or
$5-t$; the vector $d$ is determined by the line up to sign. If $d_i=d_k\ne 0$, then all cells
$q=(q_1,q_2,q_3)$ of the line satisfy $q_i=q_k$, and if $d_i=-d_k\ne 0$, they satisfy $q_i+q_k=5$. A real line
with a direction vector in $D$ contains at most six cells of the board, so distinct lines of the board lie
on distinct real lines.

\begin{proposition}[Patterns on the $6\times6\times6$ board]\label{prop:e23-min}
On the $6\times6\times6$ board:
\begin{enumerate}[label=\textup{(\alph*)},leftmargin=2.4em]
\item any four distinct lines that meet pairwise pass through one common cell;
\item every $\Ethree$ pattern has at least $13$ distinct stones, and there are $\Ethree$ patterns with
  exactly $13$;
\item there are $\Etwo$ patterns with exactly $11$ distinct stones.
\end{enumerate}
\end{proposition}

\begin{proof}
(a) Let $a_1,a_2,a_3,a_4$ be distinct lines that meet pairwise, and suppose that they do not pass through a common
cell. Distinct lines that meet lie on distinct real lines through a common point, so the direction vectors
$d^1,\dots,d^4\in D$ of the four lines are pairwise non-proportional.

We first show that the four lines lie in a plane. If no three of them pass through a common cell, then $a_1$
and $a_2$ span a plane, and each of $a_3$, $a_4$ meets $a_1$ and $a_2$ in two different cells and therefore
lies in this plane. If three of them, say $a_1,a_2,a_3$, pass through a common cell $p$, then $a_4$ does not
contain $p$ and meets $a_1,a_2,a_3$ in three different cells, which are collinear; so $a_1,a_2,a_3$ lie in the
plane spanned by $p$ and $a_4$. Hence $d^1,\dots,d^4$ lie in a two-dimensional subspace $W$ of $\mathbb{R}^3$.

We claim that $W$ contains a vector with exactly one non-zero coordinate. Suppose not. The \emph{support}
of a vector is the set of its non-zero coordinates. Let $u,v\in D\cap W$ be non-proportional vectors with
the same support. In every coordinate of the support $u_i=\pm v_i$, so exactly one of $(u+v)_i$, $(u-v)_i$
is non-zero, and it equals $\pm2$. Thus $u+v$ and $u-v$ are non-zero vectors of $W$ whose supports partition
the common support of $u$ and $v$, which has at most three elements; one of them has a support of size one, a
contradiction. So the supports of $d^1,\dots,d^4$ are four different sets of size two or three, that is, they
are $\{1,2\}$, $\{1,3\}$, $\{2,3\}$ and $\{1,2,3\}$. Let $u,v,w$ be the three of these vectors with supports
$\{1,2\}$, $\{1,3\}$, $\{1,2,3\}$. Then $w=\alpha u+\beta v$ with real $\alpha,\beta$; the second coordinate
gives $\alpha=\pm1$, the third gives $\beta=\pm1$, and then the first coordinate
$w_1=\alpha u_1+\beta v_1$ is even, a contradiction. This proves the claim.

So $W$ contains a unit vector, say the one of coordinate $k$; let $i,j$ be the other two coordinates. The
vectors of $W$ are then the vectors $z$ for which $(z_i,z_j)$ is a multiple of one fixed non-zero pair,
which can be taken from one of the vectors $d^1,\dots,d^4$ and thus, up to sign, is $(1,0)$, $(0,1)$,
$(1,1)$ or $(1,-1)$. In the first two cases $W$ is a coordinate plane, say $W=\{z: z_j=0\}$, and the vectors
of $D\cap W$ are, up to sign, the two unit vectors of the coordinates $i$ and $k$ and two vectors $u,u'$
with $u_i=u_k\ne0$ and $u'_i=-u'_k\ne0$. In the last two cases $W=\{z: z_i=\varepsilon z_j\}$ with
$\varepsilon=\pm1$, and the vectors of $D\cap W$ are, up to sign, the unit vector of coordinate $k$, one
vector with $z_k=0$, and two vectors $u,u'$ with $u_i=u_k\ne0$ and $u'_i=-u'_k\ne0$. In every case $D\cap W$
consists of exactly four vectors up to sign, so two of the four lines have the direction vectors $u$ and
$u'$. The cells $q$ of the first of them satisfy $q_i=q_k$, and those of the second satisfy $q_i+q_k=5$. A
common cell would satisfy $2q_i=5$, which is impossible. So these two lines do not meet, contrary to the
assumption.

(b) Suppose an $\Ethree$ pattern has only $12$ distinct stones. By \cref{prop:lost-stones}\,(b) the
sets $S_a\cap S_{b'}$, $S_{a'}\cap S_b$ and $S_b\cap S_{b'}$ are then non-empty, so $a$ meets $b'$, $a'$ meets
$b$, and $b$ meets $b'$. By (E3.1), $a\cap a'=\{x\}$, $a\cap b=\{c\}$ and $a'\cap b'=\{c'\}$. So the four lines
meet pairwise, and by (a) they pass through a common cell. This cell is $x$, the only common cell of $a$ and
$a'$; but $x\notin b$ by (E3.1), a contradiction. An $\Ethree$ pattern with $13$ distinct stones is the
following one, in the plane $z=0$; we write $(i,j)$ for the cell $(i,j,0)$. Lines:
$a=\{(i,0)\}$, $a'=\{(0,j)\}$, $b=\{(i,5-i)\}$, $b'=\{(i,1)\}$; cells: $x=(0,0)$, $c=(5,0)$, $c'=(0,1)$;
stone sets:
\begin{align*}
S_a&=\{(0,0),(1,0),(2,0),(3,0)\}, & S_{a'}&=\{(0,0),(0,2),(0,3),(0,5)\},\\
S_b&=\{(0,5),(1,4),(2,3),(4,1)\}, & S_{b'}&=\{(1,1),(2,1),(3,1),(4,1)\}.
\end{align*}
The conditions (E3.1) and (E3.2) are checked directly, and (E3.3) holds for $F=\emptyset$. The stone sets
meet only in $S_a\cap S_{a'}=\{(0,0)\}$, $S_{a'}\cap S_b=\{(0,5)\}$ and $S_b\cap S_{b'}=\{(4,1)\}$, so the
pattern has $16-3=13$ distinct stones.

(c) With the same notation let $l=\{(i,0)\}$, $a=\{(0,j)\}$, $b=\{(i,5-i)\}$, $x=(0,0)$, $e=(1,0)$,
$c=(0,5)$, $S_a=\{(0,0),(0,1),(0,2),(0,3)\}$ and $S_b=\{(5,0),(4,1),(3,2),(2,3)\}$. The conditions (E2.1)
and (E2.2) are checked directly, and (E2.3) holds for $F=\emptyset$. Here
$S_l=l\setminus\{e\}$ meets $S_b$ in $(5,0)$, so the pattern has $12-1=11$ distinct stones by
\cref{prop:lost-stones}\,(a).
\end{proof}

The exact distribution of the number of stones was determined by computer.

\begin{fact}[Stones of the patterns on the empty board]\label{fact:e23}
For $B=F=\emptyset$, when all $148$ lines are live, there are $5{,}152{,}320$ $\Etwo$ patterns, of which
$562{,}176$ have $11$ and $4{,}590{,}144$ have $12$ distinct stones. There are $353{,}422{,}080$ $\Ethree$
patterns, of which $10{,}948{,}608$ have $13$, $67{,}627{,}008$ have $14$, and $274{,}846{,}464$ have $15$ distinct
stones; none has $12$. Here a pattern is the tuple of its lines, cells and stone sets as listed in
\cref{def:lost}, so that an $\Ethree$ pattern and the pattern obtained from it by exchanging
$(a,b,c,S_a,S_b)$ with $(a',b',c',S_{a'},S_{b'})$ are counted separately.

\established the program \fn{e23min.cpp} in the directory \fn{logs/proof/} of the
archive (output \fn{e23min.log}; a few seconds of running time; see \cref{ssec:artifact}). The program is self-contained and constructs the
lines itself; it enumerates all patterns literally from \cref{def:lost} and counts the cells in the
union of the stone sets.
\end{fact}

Every pattern for an arbitrary pair $(B,F)$ is also a pattern for $(\emptyset,\emptyset)$, because a
non-empty $B$ only removes live lines and a non-empty $F$ only restricts the stone sets. So \cref{fact:e23} covers all patterns; it agrees with \cref{prop:lost-stones} and with
\cref{prop:e23-min}.

\smallskip
The last example of this subsection belongs to \cref{rem:conservative}(ii): it exhibits a position in
which an $\Ethree$ pattern fits although \Breaker{} has a reply that restores admissibility.

\needlines{5}
\begin{example}[A fitting pattern that does no harm]\label{ex:conservative}
In the plane $z=0$ consider the lines
\[ a=\{(i,0,0)\},\qquad a'=\{(0,j,0)\},\qquad b=\{(1,j,0)\},\qquad b'=\{(i,5-i,0)\}, \]
where $i$ and $j$ run from $0$ to $5$, the cell $x=(0,0,0)$, and the set $X'$ of the $14$ cells $(i,j,0)$
with
\begin{align*}
(i,j)\in\{&(0,0),(2,0),(3,0),(4,0),\ (0,1),(0,2),(0,3),\\
          &(1,1),(1,2),(1,3),(1,4),\ (2,3),(3,2),(4,1)\}.
\end{align*}
Let $X=X'\setminus\{x\}$, let $B$ be any set of $13$ cells with $z\ne 0$, and let $F\subseteq X$ be arbitrary;
then $\Pst(B,F,13)$ is a P-state. A line that is not contained in the plane $z=0$ has at most one cell in
this plane. The $14$ lines in the plane are live. The six lines parallel to $a$ (given by $j=0,\dots,5$) have
$4,3,3,3,1,0$ cells in $X'$, the six lines parallel to $a'$ (given by $i=0,\dots,5$) have $4,4,2,2,2,0$, the
diagonal $\{(i,i,0)\}$ has two, and $b'$ has four. So $a$, $a'$, $b$, $b'$ have four cells in $X'$ each, and
every other line has at most three. In $X$ the lines $a$ and $a'$ have three cells, and $b$, $b'$ have four
cells and meet in $(1,4,0)\in X$. So $X$ is $B$-admissible, and $X'=X+x$ has no threat. The lines with four
cells in $X'$ meet as follows: $a\cap a'=\{x\}$ and $b\cap b'=\{(1,4,0)\}$ are in $X'$,
$a'\cap b=\emptyset$, and
\[ a\cap b=\{(1,0,0)\},\qquad a\cap b'=\{(5,0,0)\},\qquad a'\cap b'=\{(0,5,0)\} \]
are empty cells. So $X'$ has exactly the danger pairs $(a,b)$, $(a,b')$ and $(a',b')$. By
\cref{lem:danger}, applied to $(a,b)$ and $(a',b')$, an $\Ethree$ pattern fits $(B,F,13)$, and indeed the
replies $(1,0,0)$ and $(0,5,0)$ leave the danger pairs $(a',b')$ and $(a,b)$, respectively. After the
reply $(5,0,0)$ only $a'$ and $b$ are live lines with four cells in $X'$, and they are disjoint: $X'$ is
$(B+(5,0,0))$-admissible. (Rule \ref{S:danger} of \cref{def:sigma} would choose $(0,5,0)$, the common
cell with the smallest id. The choice of the danger pair can matter, but only at states with $\LOST$, which a
play guided by a valid certificate never reaches.)
\end{example}

